% ==============================================================================
% Fichier : chapitres/00_introduction.tex
% Cours   : Sécurité Informatique — ENSPY, Humanités Numériques Niveau 3
% ==============================================================================

\chapter*{Introduction Générale}
\addcontentsline{toc}{chapter}{Introduction Générale}
\markboth{Introduction Générale}{Introduction Générale}

\begin{boxobjectifs}
\textbf{Objectifs de l'introduction :}
\begin{itemize}
    \item Situer le cours dans le contexte contemporain des systèmes d'information.
    \item Distinguer sécurité informatique et cybersécurité.
    \item Comprendre l'asymétrie des cyberattaques et la notion de cyberguerre.
    \item Connaître les prérequis et la structure globale du cours.
\end{itemize}
\end{boxobjectifs}

\section*{Contexte et Enjeux}

Dans le contexte contemporain des entreprises, les systèmes d'information sont devenus omniprésents, jouant un rôle essentiel dans la facilitation des tâches des employés et l'accélération des processus métier. Cette évolution résulte de l'intégration croissante d'une gamme variée d'équipements informatiques, allant des simples imprimantes aux ordinateurs complexes. Cependant, cette intégration n'est pas sans risques, car elle ouvre la voie à une interconnexion étendue entre ces équipements et les machines externes, offrant ainsi aux individus malveillants des opportunités de vol de données, de perturbation du système ou même d'arrêt complet.

Ce cours adopte une approche holistique en couvrant à la fois les concepts généraux de sécurité informatique ainsi que des éléments spécifiques pertinents pour diverses spécialisations : \textit{data scientist}, \textit{manager des systèmes d'information}, \textit{investigateur numérique}, \textit{community manager}, \textit{spécialiste en intelligence artificielle} ou \textit{expert en cybersécurité}.

\section*{Historique de la Sécurité Informatique}

\subsection*{Les origines (1940--1960)}
La sécurité informatique trouve ses racines dans les années 1940, avec les premiers ordinateurs Colossus et ENIAC. La préoccupation principale était la protection des secrets militaires et gouvernementaux, avec des techniques rudimentaires basées sur des contrôles d'accès physiques et des procédures manuelles.

\subsection*{L'essor des réseaux (1960--1980)}
L'apparition des réseaux informatiques a marqué un tournant majeur. La nécessité de sécuriser les données transmises entre les machines a donné naissance à la sécurité réseau. Les premières technologies — chiffrement, pares-feux — ont été développées durant cette période.

\subsection*{Le premier virus informatique (1988)}
En 1988, le premier virus informatique majeur, connu sous le nom de \og{}Morris\fg{}, a été créé par Robert Tappan Morris à l'Université Cornell. Ce ver a exploité une faille dans Unix pour se propager à travers ARPANET, infectant environ 6~000 ordinateurs et attirant l'attention sur la nécessité de renforcer la sécurité.

\subsection*{L'ère de la cybercriminalité (1980--2000)}
L'expansion d'Internet a ouvert la voie à une nouvelle ère de cybercriminalité : virus, malwares, attaques par déni de service (DoS). L'industrie de la sécurité informatique a connu une croissance exponentielle.

\subsection*{Évolution récente (2000--présent)}
La sécurité informatique couvre aujourd'hui réseaux, systèmes d'information, données et applications. L'intelligence artificielle, l'apprentissage automatique et l'automatisation jouent désormais un rôle crucial dans la lutte contre les cybermenaces.

\section*{Sécurité Informatique vs Cybersécurité}

\begin{boxdef}
\textbf{Cybersécurité :} Se concentre sur la protection des systèmes informatiques et des données dans le \textit{cyberespace} — appareils, réseaux, logiciels, données numériques. L'objectif est de contrer les attaques numériques : virus, malwares, intrusions, piratage.

\textbf{Sécurité Informatique :} Champ plus large. Vise à protéger l'ensemble des systèmes d'information, numériques ou non, contre les menaces physiques (vol, destruction), les erreurs humaines, les défaillances techniques et les catastrophes naturelles. Elle \emph{englobe} la cybersécurité.
\end{boxdef}

\section*{La Cyberguerre et l'Asymétrie des Attaques}

La sécurité informatique concerne aujourd'hui non seulement les entreprises, mais également les États. Leurs systèmes d'information, qui gèrent parfois des infrastructures critiques, peuvent être la cible de cyberattaques orchestrées par des acteurs étatiques ou des groupes criminels organisés.

\begin{boxalert}
\textbf{Asymétrie des cyberattaques :} Contrairement aux guerres physiques où la puissance militaire conventionnelle est déterminante, dans le cyberespace, des pays moins puissants peuvent infliger des dommages significatifs à des nations plus grandes et mieux équipées. Cette asymétrie souligne l'importance de la préparation et de la défense des infrastructures numériques.
\end{boxalert}

\section*{Nécessité et Prérequis du Cours}

La sécurité informatique est devenue une priorité absolue dans un paysage où entreprises, institutions et particuliers sont régulièrement confrontés à des cyberattaques. Ce cours de Licence~3 outille les étudiants pour :
\begin{itemize}
    \item Comprendre les principes fondamentaux de la protection des SI.
    \item Sensibiliser collègues et décideurs aux enjeux de la cybersécurité.
    \item Prévenir les attaques et garantir la résilience des SI.
\end{itemize}

\textbf{Prérequis :} Aucune connaissance préalable en sécurité informatique n'est requise. Une familiarité avec les concepts de base des réseaux informatiques et une compréhension élémentaire des algorithmes de chiffrement et de hachage constituent toutefois des atouts précieux.

\section*{Structure du Cours}

Le cours est subdivisé en neuf (9) livres :

\begin{enumerate}
    \item \textbf{Livre~I :} Définitions, concepts et bases pour le cours
    \item \textbf{Livre~II :} Security and Risk Management
    \item \textbf{Livre~III :} Sécurité des Actifs
    \item \textbf{Livre~IV :} Ingénierie et Architecture de Sécurité
    \item \textbf{Livre~V :} Réseaux et Communications
    \item \textbf{Livre~VI :} Gestion des Identités et des Accès
    \item \textbf{Livre~VII :} Évaluation et Tests de la Sécurité
    \item \textbf{Livre~VIII :} Opérations de Sécurité
    \item \textbf{Livre~IX :} Sécurité des Applications
\end{enumerate}
